\begin{lemma}[Transitivity of the lifted relation]
If $r$ is a preorder on $Q$, then for all
$X,Y,Z\in\mathsf{PowerQ}(Q)$,
\[
\mathsf{PowerRel}(r)(X,Y)\land\mathsf{PowerRel}(r)(Y,Z)
\Longrightarrow\mathsf{PowerRel}(r)(X,Z).
\].
\end{lemma}
\begin{proof}
Write $A\preceq B$ for $\mathsf{PowerRel}(r)(A,B)$ as defined in
\noderef{PowerRel}.  We prove, by well-founded induction, the assertion

\[
  A\preceq B\ \land\ B\preceq C\ \Longrightarrow\ A\preceq C
\]

on triples $(A,B,C)$, ordered lexicographically by three copies of the
child-to-parent relation \noderef{PowerChildWellFounded}.  This order is
well-founded because the child relation is well-founded.  Consequently the
induction hypothesis applies whenever one coordinate is replaced by one of
its children, with all earlier coordinates unchanged.  At each use below,
the relevant equivalence is the corresponding recursion equation for
\noderef{PowerRel}, namely one of the four reductions
\noderef{PowerRelAtomAtom}, \noderef{PowerRelAtomNode},
\noderef{PowerRelNodeAtom}, and \noderef{PowerRelNodeNode}.

First suppose that $A=\mathsf{atom}(a)$, $B=\mathsf{atom}(b)$, and
$C=\mathsf{atom}(c)$.  By \noderef{PowerRelAtomAtom}, the two hypotheses
are $r(a,b)$ and $r(b,c)$.  Transitivity of $r$, supplied by the assumed
$\mathsf{IsPreorder}$ structure, gives $r(a,c)$, and a final application of
\noderef{PowerRelAtomAtom} gives $A\preceq C$.

Now let $A=\mathsf{atom}(a)$ and $B=\mathsf{atom}(b)$.  If
$C=\mathsf{node}(\kappa,h,g)$, the second hypothesis has, by
\noderef{PowerRelAtomNode}, the form
\[
  \exists k\in\kappa,\quad \mathsf{atom}(b)\preceq g(k).
\]
Choose such a $k$.  The induction hypothesis for
$(\mathsf{atom}(a),\mathsf{atom}(b),g(k))$ is available because $g(k)$ is
a child of $C$; it combines the first hypothesis with the displayed
comparison and yields $\mathsf{atom}(a)\preceq g(k)$.  The existential
conclusion in \noderef{PowerRelAtomNode} then gives
$\mathsf{atom}(a)\preceq C$.

Next let $A=\mathsf{atom}(a)$, $B=\mathsf{node}(\iota,h,f)$, and
$C=\mathsf{atom}(c)$.  By \noderef{PowerRelAtomNode}, the first hypothesis
provides an index $i\in\iota$ with
$\mathsf{atom}(a)\preceq f(i)$.  By \noderef{PowerRelNodeAtom}, the second
hypothesis supplies $f(i)\preceq\mathsf{atom}(c)$ for every $i$, in
particular for this chosen one.  Since $f(i)$ is a child of $B$, the
induction hypothesis for $(\mathsf{atom}(a),f(i),\mathsf{atom}(c))$ gives
$\mathsf{atom}(a)\preceq\mathsf{atom}(c)$.  The atom--atom conclusion
follows through \noderef{PowerRelAtomAtom}.

If instead $C=\mathsf{node}(\kappa,h',g)$, the same first reduction gives
$i\in\iota$ with $\mathsf{atom}(a)\preceq f(i)$.  The second reduction,
\noderef{PowerRelNodeNode}, says that for this $i$ there is a
$j\in\kappa$ with $f(i)\preceq g(j)$.  Apply the induction hypothesis to
$(\mathsf{atom}(a),f(i),g(j))$: the middle entry is a child of $B$ (and the
last is a child of $C$), so the triple is lower in the induction order.
It gives $\mathsf{atom}(a)\preceq g(j)$, and the existential conclusion of
\noderef{PowerRelAtomNode} proves $\mathsf{atom}(a)\preceq C$.  This
explicitly preserves the order $\exists i\,\exists j$ arising from the
first and second comparisons.

It remains to consider a node as the leftmost term.  Let
$A=\mathsf{node}(\iota,h,f)$ and $B=\mathsf{atom}(b)$.  If
$C=\mathsf{atom}(c)$, the first hypothesis, by
\noderef{PowerRelNodeAtom}, says $f(i)\preceq\mathsf{atom}(b)$ for every
$i\in\iota$, while the second is $r(b,c)$ by
\noderef{PowerRelAtomAtom}.  For each $i$, use the induction hypothesis on
$(f(i),\mathsf{atom}(b),\mathsf{atom}(c))$; the first coordinate is a child
of $A$.  It gives $f(i)\preceq\mathsf{atom}(c)$ for every $i$, which is
exactly the universal conclusion required by \noderef{PowerRelNodeAtom}.

If $C=\mathsf{node}(\kappa,h',g)$, the second hypothesis, by
\noderef{PowerRelAtomNode}, supplies one fixed $j\in\kappa$ with
$\mathsf{atom}(b)\preceq g(j)$.  For every $i\in\iota$, the first
hypothesis gives $f(i)\preceq\mathsf{atom}(b)$ by
\noderef{PowerRelNodeAtom}.  The induction hypothesis on
$(f(i),\mathsf{atom}(b),g(j))$ is valid because $f(i)$ is a child of $A$;
it yields $f(i)\preceq g(j)$.  Thus for every $i$ we can use this same $j$,
and \noderef{PowerRelNodeNode} gives $A\preceq C$.  The quantifiers have
the required order $\forall i\,\exists j$.

Finally let both $A=\mathsf{node}(\iota,h,f)$ and
$B=\mathsf{node}(\kappa,h',g)$ be nodes.  If
$C=\mathsf{atom}(c)$, the first hypothesis and
\noderef{PowerRelNodeNode} say that for every $i\in\iota$ there is a
$j\in\kappa$ with $f(i)\preceq g(j)$.  The second hypothesis and
\noderef{PowerRelNodeAtom} say that every such $g(j)$ satisfies
$g(j)\preceq\mathsf{atom}(c)$.  For each $i$, choose the corresponding
$j$ and apply the induction hypothesis to
$(f(i),g(j),\mathsf{atom}(c))$; the first coordinate is a child of $A$.
This gives $f(i)\preceq\mathsf{atom}(c)$ for every $i$, which is the
node--atom conclusion of \noderef{PowerRelNodeAtom}.

If also $C=\mathsf{node}(\lambda,h'',k)$, then for every $i\in\iota$ the
first hypothesis, via \noderef{PowerRelNodeNode}, gives a $j\in\kappa$ with
$f(i)\preceq g(j)$.  For this $j$, the second hypothesis, again via
\noderef{PowerRelNodeNode}, gives a $\ell\in\lambda$ with
$g(j)\preceq k(\ell)$.  Apply the induction hypothesis to
$(f(i),g(j),k(\ell))$.  Its first coordinate is a child of $A$, so this is
a genuine descent even though both later coordinates also changed.  We
obtain $f(i)\preceq k(\ell)$.  Therefore, for every $i$, there exists an
$\ell$ with this comparison, and the universal-then-existential conclusion
of \noderef{PowerRelNodeNode} proves $A\preceq C$.

These eight tag triples exhaust the atom/node constructors of
$\mathsf{PowerQ}(Q)$.  In every recursive case the witnesses are chosen in
the same existential/universal order as in the four cited reductions, and
each recursive comparison is covered by the well-founded induction
hypothesis.  Hence $\mathsf{PowerRel}(r)$ is transitive.
\end{proof}
